\documentclass[11pt,twocolumn,a4paper]{article}

% ─── Packages ─────────────────────────────────────────────────────────
\usepackage[margin=2.0cm,top=2.4cm,bottom=2.4cm]{geometry}
\usepackage[utf8]{inputenc}
\usepackage[T1]{fontenc}
\usepackage{lmodern}
\usepackage{microtype}
\usepackage{amsmath,amssymb,amsthm,mathtools}
\usepackage{bm}
\usepackage{graphicx}
\usepackage{booktabs}
\usepackage{array}
\usepackage{enumitem}
\usepackage{float}
\usepackage{caption}
\usepackage{natbib}
\usepackage{xcolor}
\usepackage{listings}
\usepackage[colorlinks=true,linkcolor=blue!60!black,
            citecolor=blue!60!black,urlcolor=blue!60!black]{hyperref}
\usepackage{fancyhdr}
\usepackage{titlesec}
\usepackage{cleveref}

% ─── Page Style ───────────────────────────────────────────────────────
\pagestyle{fancy}
\fancyhf{}
\fancyhead[L]{\small\textit{Sangoma: Declaring an Instrument by its Target}}
\fancyhead[R]{\small\thepage}
\renewcommand{\headrulewidth}{0.4pt}

\titleformat{\section}{\large\bfseries}{\thesection.}{0.5em}{}
\titleformat{\subsection}{\normalsize\bfseries}{\thesubsection.}{0.5em}{}

% ─── Theorem Environments ─────────────────────────────────────────────
\newtheorem{theorem}{Theorem}[section]
\newtheorem{lemma}[theorem]{Lemma}
\newtheorem{corollary}[theorem]{Corollary}
\newtheorem{proposition}[theorem]{Proposition}
\theoremstyle{definition}
\newtheorem{definition}[theorem]{Definition}
\newtheorem{example}[theorem]{Example}
\newtheorem{principle}[theorem]{Principle}
\theoremstyle{remark}
\newtheorem{remark}[theorem]{Remark}

% ─── Custom Commands ──────────────────────────────────────────────────
\newcommand{\R}{\mathbb{R}}
\newcommand{\N}{\mathbb{N}}
\newcommand{\Rpos}{\R_{>0}}
\newcommand{\flo}{\beta}
\newcommand{\thick}{t}
\newcommand{\pot}{\kappa}
\newcommand{\Om}{\Omega}
\newcommand{\Gph}{\Gamma}
\newcommand{\wt}{w}
\newcommand{\cut}{\operatorname{cut}}
\newcommand{\res}{\varrho}
\newcommand{\med}{\mathcal{M}}

% ─── Listings ─────────────────────────────────────────────────────────
\lstdefinestyle{sgn}{
  basicstyle=\ttfamily\footnotesize,
  keywordstyle=\bfseries,
  commentstyle=\itshape\color{gray!80!black},
  showstringspaces=false,
  breaklines=true,
  columns=fullflexible,
  frame=single,
  framesep=4pt,
  rulecolor=\color{gray!50},
  morekeywords={floor,medium,species,source,motion,construct,stage,chain,
                assemble,from,target,via,rung,at,observe,assert,emit,let,as,
                ceiling,expresses,couples,by}
}
\lstset{style=sgn}

\captionsetup{font=small,labelfont=bf}

% ══════════════════════════════════════════════════════════════════════
\begin{document}

\twocolumn[
\begin{@twocolumnfalse}
\begin{center}

{\LARGE\bfseries Sangoma: Declaring an Instrument by its Target\\[4pt]
A Constructor Language in which Unreachability is a Diagnostic,\\
not a Rendering Failure}

\vspace{12pt}

{\large Kundai Farai Sachikonye}

\vspace{6pt}

{\small\itshape Independent Researcher \\[0.3em]
\texttt{kundai.sachikonye@wzw.tum.de}}

\vspace{6pt}{\small August 2026}

\vspace{16pt}

\begin{minipage}{0.92\textwidth}
\textbf{Abstract.}
We specify \emph{sangoma}, a language for building sounds by declaring what the
finished sound must satisfy rather than by naming the steps that produce it. The
argument for the target form is that a procedure both over-determines and
under-determines: many chains reach the same sound, so naming one excludes the
rest, and a chain that cannot reach its intended sound gives no warning until it
has been rendered. A sangoma program instead states the properties required and
the resolution at which each is asserted, and the compiler answers a question no
procedure poses---\emph{is this reachable at this resolution?} If it is, a chain
is emitted; if not, the program is refused before rendering and the refusal names
the depth that would be required, $\lceil\log(1-\pot^{\star})/\log(1-\pot)\rceil$
rungs at the strongest available stage. We prove that a single processing stage
cannot be refined without bound, since the cost of sharpening a separator
diverges while its return is bounded by the extent of the space, and that fine
resolution is therefore reached by composing ordinary stages in series; the
crossover is measured, with composition overtaking sharpening at a target of
$0.95$ and exceeding it twentyfold by $0.999$. We give a type system carrying
resolution and separator thickness in the type, in which a target stating no
resolution is rejected, a target finer than the declared ambient floor is
rejected as an over-claim, and a stage whose separator engulfs the assembly is
statically unobservable and therefore refused rather than rendered and found
silent. Physical constraints---headroom against a declared ceiling, latency
accumulation, and negotiation between a program's floor and the render path's
numeric resolution---are enforced as rules, while everything above them is
composable and unjudged, because a language that refused on taste would enforce
its author's present opinions. We give the complete grammar of the \texttt{.sgn}
file, a worked program, a prototype compiler, and thirteen deterministic
experiments. Nothing outside this paper is required to read it.
\end{minipage}

\vspace{8pt}
\noindent\textbf{Keywords:}
domain-specific language, sound design, constructor calculus, reachability,
separator, detectability, headroom, type system, static diagnostics

\vspace{18pt}
\end{center}
\end{@twocolumnfalse}
]

%──────────────────────────────────────────────────────────────────────
\section{Introduction}
\label{sec:intro}
%──────────────────────────────────────────────────────────────────────

\subsection{Why a target and not a procedure}

A producer building a bass sound has, at any moment, a great many plugins
available and no way to know in advance whether the chain they are assembling can
reach the sound they want. They find out by rendering. If it does not, they
adjust and render again, and the loop is the ordinary experience of sound design.

We take the position that this loop is not intrinsic to the work but an artefact
of describing chains procedurally. A procedure is a poor specification for two
reasons, and they pull in opposite directions.

It \emph{over-determines}. Many chains reach the same sound; naming one excludes
the rest, including the ones that are cheaper, shorter, or more robust to the
material changing. A specification that names a procedure has committed to a
route before the destination has been stated.

It \emph{under-determines}. A chain that cannot reach the intended sound is
syntactically indistinguishable from one that can. Nothing in the description
carries the information that would settle it, so nothing can warn, and the
producer learns by rendering.

Sangoma therefore declares a \emph{target}: what the finished sound must satisfy,
and the resolution at which each requirement is asserted. The compiler then
decides reachability (\Cref{prop:reach}), emits a chain when the target is
reachable, and refuses with the required depth named when it is not. The loop
above collapses into a compile-time diagnostic.

\subsection{What the language will and will not judge}

A language that refuses programs must be careful about the grounds on which it
refuses, and this is the design decision we most want to state plainly.

Some constraints in audio are physical. A chain whose summed gain drives the
signal past a declared ceiling will clip; a chain of look-ahead stages incurs
latency that must be compensated; a claim to resolve a difference finer than the
render path can represent is a claim about an artefact. These are checkable and
sangoma checks them (\S\ref{sec:physical}).

Almost everything else in sound design is taste. Whether a compressor belongs
before or after a saturator, whether a resample should precede a filter---these
have no correct answer, and a language that refused on such grounds would encode
its author's present opinions and then enforce them. Worse, it would obstruct
exactly the misuse that produces new sounds: the genres this language is built
for were made substantially by people using tools wrongly on purpose.

Sangoma therefore refuses on physics and on unreachability, and is silent above
that. The compiler will tell a producer that a chain cannot reach a stated
target, or that it will clip, or that it claims a resolution nothing delivers. It
will not tell them a chain is a bad idea.

\subsection{Contributions}

\begin{enumerate}[label=(\roman*),leftmargin=2.2em]
\item A theory of processing stages as separators: thickness, additivity, and the
  detectability criterion (\S\ref{sec:separators}).
\item The divergence of refinement cost, the consequent necessity of composing
  stages, and a measured crossover (\S\ref{sec:composition}).
\item The reachability check, which decides a target statically and reports the
  required depth on failure (\S\ref{sec:reachability}).
\item A constructor calculus of species, stages and chains
  (\S\ref{sec:calculus}).
\item The grammar of \texttt{.sgn}, and a type system carrying resolution and
  thickness (\S\ref{sec:grammar}).
\item Physical rules: headroom, latency, and floor negotiation
  (\S\ref{sec:physical}).
\item A worked program, a prototype compiler, and thirteen experiments
  (\S\ref{sec:worked}--\S\ref{sec:validation}).
\end{enumerate}

%──────────────────────────────────────────────────────────────────────
\section{Substrate}
\label{sec:substrate}
%──────────────────────────────────────────────────────────────────────

We restate, self-containedly, the structure the language is defined over.

\begin{definition}[Contact graph, medium, cut, residue, floor]
\label{def:substrate}
A \emph{contact graph} is $\Gph=(V,E,\wt)$ with $V$ finite and non-empty, $E$ a
set of unordered pairs, and $\wt\colon E\to\Rpos$; terminology follows
\citet{diestel2010graph}. A distinguished vertex $\med$, the \emph{medium}, is
adjacent to every other vertex. For $S$ containing an item and omitting $\med$,
the \emph{cut} is $\cut(S)=\{\{u,v\}\in E: u\in S,\ v\notin S\}$ and the
\emph{residue} is $\res(S)=\sum_{e\in\cut(S)}\wt(e)$. The \emph{floor} is
$\flo=\min_{e\in E}\wt(e)$ and the \emph{total} is $\Om=\sum_{e\in E}\wt(e)$.
\end{definition}

\begin{lemma}[Positive floor]
\label{lem:floor}
$\flo>0$ and every cut satisfies $\res(S)\ge\flo>0$.
\end{lemma}

\begin{proof}
$\flo$ is the minimum of a finite non-empty set of positive reals. Since $S$
contains an item and omits $\med$, and every item is adjacent to $\med$, the cut
is non-empty and its weight is a sum of terms each at least $\flo$.
\end{proof}

\begin{remark}[The floor in a render path]
\label{rem:audio-floor}
The floor is an ordinary engineering quantity. A converter of $b$ bits over a
range $D$ does not distinguish differences below $D/2^{b}$; a spectral analysis of
window length $N$ at sample rate $f_s$ does not resolve frequencies closer than
$f_s/N$ \citep{harris1978use}. These fix the weights of \Cref{def:substrate} in
practice, and they are read from specifications rather than estimated.
\end{remark}

%──────────────────────────────────────────────────────────────────────
\section{Stages as Separators}
\label{sec:separators}
%──────────────────────────────────────────────────────────────────────

\begin{definition}[Stage, thickness]
\label{def:stage}
A \emph{stage} is a processing operation together with its \emph{thickness}
$\thick>0$: the resolution at which the stage's effect can be told from its
absence. A stage of thickness $\thick$ makes no distinction finer than $\thick$.
\end{definition}

\begin{principle}[A stage must report]
\label{prin:report}
A stage is declared by what it does \emph{and} by what can be measured about what
it did. A stage whose effect admits no measurement is a declaration error, not a
permissible design.
\end{principle}

The principle has a practical reading. A producer who inserts a plugin and cannot
tell, by any measurement or by ear, whether it is engaged has not added a stage
to the chain; they have added an unfalsifiable belief about the chain.
\Cref{thm:detect} makes the condition precise.

\begin{proposition}[Thickness is additive along a chain]
\label{prop:additive}
A chain of stages of thicknesses $\thick_1,\dots,\thick_n$ has thickness
$\sum_i \thick_i$, and no downstream stage is sharper than the chain that
precedes it.
\end{proposition}

\begin{proof}
Each stage introduces an uncertainty of $\thick_i$ in the attribution of the
chain's output to its input, and these accumulate additively to first order since
the stages act in series on the same signal. The sum is non-decreasing in $n$, so
the chain's resolution never improves by the addition of a stage.
\end{proof}

\begin{remark}[Why nothing gets sharper downstream]
\label{rem:no-sharpening}
\Cref{prop:additive} is the formal content of a familiar studio fact: a
measurement taken at the end of a chain cannot attribute a difference to a stage
more precisely than the intervening stages permit. A chain is an instrument, and
an instrument's resolution is set by all of it.
\end{remark}

The experiment \texttt{sangoma\_thickness\_additive} confirms additivity over a
three-stage chain and that the running total never decreases.

\begin{definition}[Detectability]
\label{def:detect}
Let an assembly have extent $\Lambda>0$---the range over which its behaviour
varies. The \emph{detectability ratio} of a stage of thickness $\thick$ within
that assembly is $\tau = \thick/\Lambda$.
\end{definition}

\begin{theorem}[Detectability criterion]
\label{thm:detect}
A stage is observable within an assembly if and only if $\tau<1$. A stage with
$\tau\ge1$ is engulfed by its own separator: no measurement of the assembly
distinguishes its presence from its absence.
\end{theorem}

\begin{proof}
The assembly's behaviour varies over a range $\Lambda$. A stage of thickness
$\thick$ makes no distinction finer than $\thick$ by \Cref{def:stage}. If
$\thick\ge\Lambda$ then the whole range over which the assembly varies lies
inside one indistinguishable band of the stage, so every configuration of the
assembly falls in the same band and no measurement separates them. If
$\thick<\Lambda$ there exist two configurations differing by more than $\thick$,
which the stage distinguishes.
\end{proof}

\begin{remark}[Refused rather than rendered and found silent]
\label{rem:refuse-silent}
\Cref{thm:detect} is enforced by the type system (rule \textsc{(Obs)},
\S\ref{sec:grammar}), so an engulfed stage is rejected statically. This is the
difference between a compiler saying ``this stage cannot have an audible effect
in this chain'' and a producer rendering, hearing nothing, and spending an hour
deciding whether the plugin is broken.
\end{remark}

The experiment \texttt{sangoma\_detectability} sweeps $\tau$ across
$\{0.1,0.5,1.0,2.0\}$ and finds the first two observable and the last two not.

%──────────────────────────────────────────────────────────────────────
\section{Why Chains and not Sharper Stages}
\label{sec:composition}
%──────────────────────────────────────────────────────────────────────

\begin{definition}[Refinement cost]
\label{def:refinecost}
Let $\Gamma(\thick)$ be the cost of realising a separator of thickness $\thick>0$,
non-increasing in $\thick$ with $\Gamma(\thick)\to\infty$ as $\thick\to0^{+}$.
\end{definition}

\begin{theorem}[Refinement halts above zero]
\label{thm:refine}
With return bounded above by the extent $\Lambda$, marginal return per unit cost
tends to zero as $\thick\to0^{+}$, and construction halts at a strictly positive
thickness.
\end{theorem}

\begin{proof}
Return is bounded by $\Lambda$ regardless of $\thick$, so the return available
from further thinning is non-increasing in the thinning already done, while cost
diverges by \Cref{def:refinecost}. A bounded non-increasing numerator over a
divergent denominator tends to zero and falls below any positive threshold at
some $\thick^{\ast}>0$.
\end{proof}

\begin{definition}[Rung, power, ladder]
\label{def:rung}
A \emph{rung} is a stage considered as closing part of the distance to a target;
its \emph{power} $\pot\in[0,1]$ is the fraction of the above-floor gap it closes.
A \emph{ladder} is a series composition of rungs.
\end{definition}

\begin{theorem}[Composition]
\label{thm:compose}
A ladder of powers $\pot_1,\dots,\pot_n$ attains composite power
\[
  \pot_{\mathrm{total}} = 1-\prod_{i=1}^{n}(1-\pot_i).
\]
\end{theorem}

\begin{proof}
Each rung leaves a fraction $1-\pot_i$ of the above-floor gap unclosed; series
composition multiplies the surviving fractions.
\end{proof}

\begin{corollary}[Required depth]
\label{cor:rungcount}
Reaching composite power $\pot^{\star}<1$ with rungs of power $\pot$ requires
$n^{\star}=\lceil \log(1-\pot^{\star})/\log(1-\pot)\rceil$ rungs.
\end{corollary}

\begin{proof}
Solve $1-(1-\pot)^{n}\ge\pot^{\star}$ for $n$.
\end{proof}

\begin{table}[H]
\centering
\small
\caption{Composition against sharpening, per-rung power $\pot=0.35$ and cost
$\Gamma(\thick)=1/\thick$ (\texttt{sangoma\_nesting\_beats\_sharpening}).}
\label{tab:crossover}
\begin{tabular}{@{}rrrrr@{}}
\toprule
Target & Single & Rungs & Ladder & Ratio \\
\midrule
$0.90$   & $10$   & $6$  & $17.1$ & $0.58$ \\
$0.95$   & $20$   & $7$  & $20.0$ & $1.00$ \\
$0.99$   & $100$  & $11$ & $31.4$ & $3.18$ \\
$0.999$  & $1000$ & $17$ & $48.6$ & $20.6$ \\
$0.9999$ & $10^4$ & $23$ & $65.7$ & $152$ \\
\bottomrule
\end{tabular}
\end{table}

\begin{remark}[The crossover is honest about where it sits]
\label{rem:crossover}
\Cref{tab:crossover} shows composition losing at a loose target, tying exactly at
$0.95$, and winning by two orders of magnitude at a tight one. We report the
crossover rather than only the regime that flatters the design: a single stage is
genuinely the right answer for a coarse target, and the case for chains is
asymptotic. That it is asymptotic is not a weakness, because the targets that
matter in sound design---a bass that survives a club system, a master that holds
up against a reference---are the tight ones.
\end{remark}

\begin{remark}[What a resample chain is]
\label{rem:resample}
The practice of bouncing a sound, processing the bounce, bouncing again, and
repeating---resampling, in the vernacular---is exactly a ladder in the sense of
\Cref{def:rung}. Each pass is an ordinary stage of moderate power; the composite
reaches a resolution no single stage of comparable cost reaches. This is why the
genres built on that practice sound the way they do, and \Cref{thm:refine} is why
no amount of care applied to a single plugin substitutes for it.
\end{remark}

%──────────────────────────────────────────────────────────────────────
\section{Reachability}
\label{sec:reachability}
%──────────────────────────────────────────────────────────────────────

\begin{definition}[Target]
\label{def:target}
A \emph{target} is a finite set of requirements, each naming a measurable
property, a relation, a magnitude, and the resolution at which the requirement is
asserted. A requirement stating no resolution is ill-formed.
\end{definition}

\begin{proposition}[Reachability is decidable statically]
\label{prop:reach}
Given a target requiring composite power $\pot^{\star}$ and available stages of
powers $\pot_1,\dots,\pot_m$, reachability is decided by \Cref{thm:compose}. If
the best attainable composite is below $\pot^{\star}$, the program is rejected and
the refusal reports the depth
\[
  n^{\star} = \left\lceil
    \frac{\log(1-\pot^{\star})}{\log(1-\pot_{\max})}
  \right\rceil
\]
required with the strongest available stage.
\end{proposition}

\begin{proof}
Composite power is a computable function of the declared stage powers
(\Cref{thm:compose}), so the comparison is decidable before rendering. The depth
follows from \Cref{cor:rungcount} with $\pot=\pot_{\max}$, which minimises $n$
since $n$ decreases in $\pot$.
\end{proof}

\begin{remark}[Unreachability is a diagnostic]
\label{rem:diag}
\Cref{prop:reach} is what distinguishes a target declaration from a procedural
one. A chain that cannot reach its target is refused before any audio is
rendered, and the refusal carries the required depth: the producer learns not
merely that the target is out of reach but by how much.
\end{remark}

The experiment \texttt{sangoma\_reachability\_reports\_depth} confirms that the
reported depth agrees with the closed form across three targets, and
\texttt{sangoma\_reachability\_diagnostic} confirms that a weak chain is refused
while a sufficient one is not.

%──────────────────────────────────────────────────────────────────────
\section{The Constructor Calculus}
\label{sec:calculus}
%──────────────────────────────────────────────────────────────────────

\begin{definition}[Species]
\label{def:species}
A \emph{species} declares a sound-generating element by its source and its
motion: the parameters fixing what it emits at rest, and those fixing how that
emission varies. A species is not a preset; it is the smallest object the
language individuates.
\end{definition}

\begin{definition}[Chain, assembly]
\label{def:chain}
A \emph{chain} is a series composition of stages applied to a species. An
\emph{assembly} is a chain considered together with the target it is required to
satisfy.
\end{definition}

\begin{remark}[Why the calculus is small]
\label{rem:small}
There is one composition rule---series application of a stage---and no
user-defined functions, lambdas, loops or recursion. The only iteration is the
implicit one inside a ladder, whose length the compiler derives from the target
by \Cref{cor:rungcount}. This is a deliberate restriction: a language rich enough
to express arbitrary computation over chains would be rich enough to defeat the
static reachability check, which is the one thing the language exists to provide.
\end{remark}

%──────────────────────────────────────────────────────────────────────
\section{Grammar and Type System}
\label{sec:grammar}
%──────────────────────────────────────────────────────────────────────

\subsection{Lexical structure}

Source is UTF-8 with extension \texttt{.sgn}. Comments run from \texttt{--} to
end of line; the language is not layout-sensitive. Identifiers are
\texttt{[A-Za-z\_][A-Za-z0-9\_]*}; numbers are
\texttt{-?[0-9]+(.[0-9]+)?([eE][+-]?[0-9]+)?}.

Keywords: \texttt{floor}, \texttt{medium}, \texttt{species}, \texttt{source},
\texttt{motion}, \texttt{construct}, \texttt{stage}, \texttt{chain},
\texttt{assemble}, \texttt{from}, \texttt{target}, \texttt{via}, \texttt{rung},
\texttt{at}, \texttt{observe}, \texttt{assert}, \texttt{emit}, \texttt{let},
\texttt{as}, \texttt{ceiling}, \texttt{expresses}, \texttt{couples},
\texttt{by}.

A numeric literal may carry an explicit resolution by the suffix
\texttt{value\#floor}. A literal without one inherits the ambient floor, and
\emph{no literal of zero resolution is writable}: the lexer rejects a
non-positive floor suffix.

\subsection{Context-free grammar}

\begin{lstlisting}[caption={Sangoma grammar (EBNF).},label={lst:ebnf}]
program    ::= { decl }
decl       ::= floorDecl | mediumDecl | speciesDecl
             | construct | stmt
floorDecl  ::= "floor" num
mediumDecl ::= "medium" ident "{" { field } "}"
speciesDecl::= "species" ident "{"
                 "source" "{" { field } "}"
                 "motion" "{" { field } "}"
               "}"

construct  ::= "construct" ident "{"
                 { compose }
                 "target" "{" { targetItem } "}"
               [ "via" "{" ladder "}" ]
               "}"
compose    ::= "stage" ident [ ":=" expr ]
             | "assemble" ident "from" identList
targetItem ::= ident relop num
ladder     ::= rungExpr { ">>" rungExpr }
rungExpr   ::= "rung" ident [ "at" num ]

stmt       ::= "observe" expr [ "as" ident ]
             | "assert" cond [ "emit" str ]
             | bind
bind       ::= [ "let" ] ident ":=" expr
field      ::= ident ":" (num | str)
expr       ::= ident | num | str | call | "(" expr ")"
call       ::= ident "(" [ argList ] ")"
argList    ::= arg { "," arg }
arg        ::= [ ident ":" ] expr
cond       ::= expr relop expr
relop      ::= "<" | ">" | "<=" | ">=" | "==" | "!="
identList  ::= ident { "," ident }
num        ::= NUM [ "#" NUM ]
\end{lstlisting}

\subsection{Types}

\begin{definition}[Types and judgement]
\label{def:types}
\[
  \tau ::= \textsf{Species} \mid \textsf{Stage} \mid \textsf{Assembly}
        \mid \textsf{Scalar} \mid \textsf{Target} \mid \textsf{Plan}
\]
The judgement, in the style of \citet{pierce2002types} with soundness after
\citet{wright1994syntactic}, is $\Gamma\vdash e:\tau\,@\,\flo\,/\,\thick$: the
expression $e$ has type $\tau$, resolution $\flo>0$, and separator thickness
$\thick\ge\flo$.
\end{definition}

\[
\frac{\Gamma \vdash n \quad \flo > 0}
     {\Gamma \vdash n\#\flo : \textsf{Scalar} \,@\, \flo \,/\, \flo}
\ \textsc{(Lit)}
\]

\[
\frac{\Gamma \vdash s : \textsf{Species} \,@\, \flo \,/\, \thick_s}
     {\Gamma \vdash \texttt{stage}(s) : \textsf{Stage} \,@\, \flo \,/\, \thick_s}
\ \textsc{(Stg)}
\]

\[
\frac{\Gamma \vdash c_1, c_2 : \textsf{Stage} \,@\, \flo \,/\, \thick_i}
     {\Gamma \vdash \texttt{chain}(c_1,c_2) :
      \textsf{Assembly} \,@\, \flo \,/\, \thick_1 + \thick_2}
\ \textsc{(Chn)}
\]

\[
\frac{\Gamma \vdash A : \textsf{Assembly} \,@\, \flo \,/\, \thick \quad
      \tau(A) < 1}
     {\Gamma \vdash \texttt{observe}(A) :
      \textsf{Assembly} \,@\, \flo \,/\, \thick}
\ \textsc{(Obs)}
\]

Rule \textsc{(Chn)} is \Cref{prop:additive} at the type level; rule
\textsc{(Obs)} is \Cref{thm:detect}, so an engulfed assembly is rejected
statically rather than measured and found silent.

\begin{theorem}[No zero-resolution value]
\label{thm:no-zero}
There is no well-typed sangoma expression of resolution $0$; a claim of exact
measurement is not expressible.
\end{theorem}

\begin{proof}
Every value-introducing rule carries the side condition $\flo>0$, and no rule
lowers a resolution annotation: \textsc{(Stg)} and \textsc{(Obs)} carry $\flo$
through and \textsc{(Chn)} increases the thickness. Induction on the derivation
gives the result.
\end{proof}

\subsection{Static diagnostics}

\begin{proposition}[Resolution required on a target]
\label{prop:target-res}
A target requirement whose magnitude carries no resolution annotation is
rejected at parse time.
\end{proposition}

\begin{proof}
By \Cref{def:target} a requirement asserts a property at a resolution; a
magnitude without one asserts a distinction of unstated width, which by
\Cref{thm:no-zero} the language does not express.
\end{proof}

\begin{proposition}[Over-claim]
\label{prop:overclaim}
A target asserting a resolution finer than the ambient floor is rejected, with
both quantities named.
\end{proposition}

\begin{proof}
The ambient floor is the finest distinction the declared instruments support
(\Cref{lem:floor}); a requirement asserting a finer one asserts a distinction no
instrument delivers, and the comparison is decidable statically.
\end{proof}

\begin{principle}[Every refusal carries a remedy]
\label{prin:remedy}
Each diagnostic names the rule that fired, what is wrong, and what to do about
it. The checker emits no refusal without a way forward.
\end{principle}

The experiments \texttt{sangoma\_target\_needs\_resolution},
\texttt{sangoma\_over\_claim\_refused} and
\texttt{sangoma\_diagnostics\_carry\_remedy} confirm these; the last checks that
every diagnostic produced across three malformed programs carries a non-empty
remedy.

%──────────────────────────────────────────────────────────────────────
\section{Physical Rules}
\label{sec:physical}
%──────────────────────────────────────────────────────────────────────

These are the constraints the language enforces, and they are enforced because
they are facts about signals rather than opinions about music.

\begin{proposition}[Headroom]
\label{prop:headroom}
A chain whose declared stage gains sum to a peak above the declared ceiling is
refused.
\end{proposition}

\begin{proof}
Peak level after a series of gain stages is the sum of the gains in decibels
applied to the input peak. If that exceeds the ceiling the signal is not
representable in the output format and will clip
\citep{lund2006stop}.
\end{proof}

\begin{proposition}[Latency]
\label{prop:latency}
The latency of a chain is the sum of its stages' declared latencies, and must be
reported so that it can be compensated.
\end{proposition}

\begin{proof}
Stages act in series, so their delays compose additively.
\end{proof}

\begin{definition}[Numeric resolution]
\label{def:numres}
The \emph{numeric resolution} $\varepsilon_T$ of a render path is the greatest
absolute error it may introduce: accumulated floating-point rounding
\citep{goldberg1991what} plus the quantisation step of the output format.
\end{definition}

\begin{theorem}[Floor negotiation]
\label{thm:negotiate}
A program of ambient floor $\flo$ may render on a path of numeric resolution
$\varepsilon_T$, against a reference of resolution $\varepsilon_R$, only if
$\varepsilon_R+\varepsilon_T\le\flo$. Otherwise the distinctions it asserts are
not separable from the errors it incurs, and it must be refused before rendering.
\end{theorem}

\begin{proof}
Two magnitudes differing by less than $\varepsilon_R+\varepsilon_T$ are not
distinguishable by computation on that path, since each may be perturbed by up to
its resolution. The floor asserts that differences of size $\flo$ are meaningful;
if $\varepsilon_R+\varepsilon_T>\flo$ then a difference the program treats as
meaningful lies inside the machinery's error, and a reported distinction may be
an artefact.
\end{proof}

\begin{remark}[The two quantities are independent]
\label{rem:indep}
It is tempting to suppose the render path's error is automatically bounded by the
floor. It is not: the floor is a property of the weights the author declared,
$\varepsilon_T$ a property of the arithmetic and output format. Nothing couples
them, which is why the check must be performed rather than assumed.
\end{remark}

The experiments \texttt{sangoma\_headroom\_refusal},
\texttt{sangoma\_latency\_accounted} and \texttt{sangoma\_floor\_negotiation}
confirm these, the last checking that the refusal names both the program's floor
and the path's resolution.

\begin{remark}[Where the rules stop]
\label{rem:rules-stop}
The list above is complete. There is no rule about stage ordering, none about
which processor suits which material, and none that prefers one topology to
another. \S\ref{sec:intro} gives the reason: those questions have no answer the
language is entitled to enforce, and a compiler that answered them would be
obstructing the misuse that the idiom is largely made of.
\end{remark}

%──────────────────────────────────────────────────────────────────────
\section{Worked Program}
\label{sec:worked}
%──────────────────────────────────────────────────────────────────────

\begin{lstlisting}[caption={A sangoma program constructing a bass.},label={lst:prog}]
-- reese.sgn
floor 0.02

medium air {
  ceiling: -1.0
}

species reese {
  source { operators: 2, ratio: 1.0, index: 3.5 }
  motion { rate: 0.3, depth: 0.8 }
}

construct bass {
  stage fm_source
  stage resample
  stage saturate

  target {
    crest     >= 6.0#0.5
    midrange  >= 0.40#0.05
    width     <= 0.85#0.02
  }

  via { rung fm_source at 0.40
     >> rung resample  at 0.35
     >> rung saturate  at 0.55 }
}

observe bass as result
assert result.tau < 1.0
  emit "a stage is engulfed by its own separator"
\end{lstlisting}

The ladder attains composite power
$1-(1-0.40)(1-0.35)(1-0.55)=1-0.60\times0.65\times0.45=0.8245$. Three stages of
differing power are used rather than three of one, per \Cref{rem:crossover}. The
final assertion writes out the detectability check of \Cref{thm:detect}, which
rule \textsc{(Obs)} enforces regardless.

The experiment \texttt{sangoma\_parses\_worked\_example} recovers from this
source a floor of $0.02$, two declared species, three stages, three target
requirements, and a composite power of $0.8245$.

\begin{remark}[Reading the targets]
\label{rem:targets}
Each requirement carries the resolution at which it is asserted, and the
resolutions are not decorative. \texttt{crest >= 6.0\#0.5} asserts a crest factor
of at least $6$\,dB, resolved to half a decibel---which is what an ordinary meter
delivers. Writing \texttt{6.0\#0.001} instead would be refused by
\Cref{prop:overclaim}, because no instrument in the declared chain resolves a
thousandth of a decibel and the program would be asserting a distinction that
does not exist.
\end{remark}

%──────────────────────────────────────────────────────────────────────
\section{Prototype and Validation}
\label{sec:validation}
%──────────────────────────────────────────────────────────────────────

A prototype compiler accompanies this paper---lexer, recursive-descent parser,
type checker and static diagnostics---in \texttt{hkcore.py}, with experiments in
\texttt{exp\_sangoma.py}. It uses the Python standard library only. The suite
across all three papers reports \textbf{40 run, 40 passed, 0 failed} in
$0.16$\,s, of which thirteen experiments belong to this paper; every seed is
fixed and re-running reproduces each record bit-identically.

\begin{table}[H]
\centering
\small
\caption{Experiments for this paper.}
\label{tab:experiments}
\begin{tabular}{@{}lp{4.2cm}@{}}
\toprule
Experiment & Establishes \\
\midrule
\texttt{parses\_worked\_example} & \Cref{lst:prog}; $\pot=0.8245$ \\
\texttt{target\_needs\_resolution} & \Cref{prop:target-res} at parse time \\
\texttt{over\_claim\_refused} & \Cref{prop:overclaim} with remedy \\
\texttt{reachability\_reports\_depth} & \Cref{prop:reach}; 3 targets agree \\
\texttt{reachability\_diagnostic} & weak chain refused, strong accepted \\
\texttt{detectability} & \Cref{thm:detect}; $\tau$ swept over $4$ values \\
\texttt{thickness\_additive} & \Cref{prop:additive}; total $0.10$ \\
\texttt{headroom\_refusal} & \Cref{prop:headroom} \\
\texttt{latency\_accounted} & \Cref{prop:latency}; $2560$ samples \\
\texttt{floor\_negotiation} & \Cref{thm:negotiate}; both figures named \\
\texttt{nesting\_beats\_sharpening} & \Cref{tab:crossover}; crossover at $0.95$ \\
\texttt{target\_required} & a construct with no target is refused \\
\texttt{diagnostics\_carry\_remedy} & \Cref{prin:remedy}; all non-empty \\
\bottomrule
\end{tabular}
\end{table}

\begin{remark}[A claim withdrawn]
\label{rem:withdrawn}
\texttt{sangoma\_nesting\_beats\_sharpening} was first written to assert that a
ladder is cheaper than a single sharp stage at a target of $0.95$. Under the
reciprocal cost law the two are exactly equal there, and the equality is an
artefact of that cost law rather than a result. The experiment was rewritten to
report the crossover table, and \Cref{rem:crossover} states the weaker asymptotic
claim that survives. We record this because the original claim would have been
easy to state and impossible to defend.
\end{remark}

%──────────────────────────────────────────────────────────────────────
\section{Related Notions}
\label{sec:related}
%──────────────────────────────────────────────────────────────────────

\emph{Unit-generator languages.} Music-DSP environments
\citep{puckette1996pure,mccartney2002rethinking} describe signal flow
procedurally: a patch names its processors and their connections. Sangoma names
neither, and derives a chain from a target instead; the reachability check of
\Cref{prop:reach} has no analogue in a patch language, because a patch carries no
statement of what it is trying to achieve.

\emph{Synthesis by analysis and by optimisation.} Automatic synthesiser
programming \citep{horner1993machine} searches parameter space for a match to a
recorded target. The difference is in what is declared: such a system takes an
audio exemplar and returns parameters, while sangoma takes stated properties with
stated resolutions and returns a chain or a refusal. Neither subsumes the other,
and a producer with an exemplar should use the former.

\emph{Constraint and target-driven design.} Declaring an object by the properties
it must satisfy is the ordinary form in constraint programming
\citep{rossi2006handbook}. What is added here is resolution: every requirement
carries the width at which it is asserted, and \Cref{prop:overclaim} refuses a
constraint finer than the instruments support.

\emph{Loudness and headroom practice.} The physical rules of
\S\ref{sec:physical} formalise established mastering practice
\citep{lund2006stop} rather than proposing anything new; the contribution is that
they are checked before rendering rather than discovered after.

%──────────────────────────────────────────────────────────────────────
\section{Limits}
\label{sec:limits}
%──────────────────────────────────────────────────────────────────────

\emph{Stage powers must be calibrated.} \Cref{def:rung} takes a stage's power as
declared. Establishing what a particular saturator actually closes, on particular
material, is a measurement problem this paper does not solve. A miscalibrated
registry yields a well-typed program that reaches the wrong sound, and the
compiler cannot detect this.

\emph{Additivity is first-order.} \Cref{prop:additive} treats thicknesses as
summing, which holds when stages act approximately independently. Strongly
interacting stages---a compressor feeding a distortion whose drive changes the
compressor's own detection---violate it. We regard this as the most serious open
problem in the design, and note that the interacting case is common in exactly
the chains the language is meant for.

\emph{No claim about audibility.} Detectability (\Cref{thm:detect}) is a
statement about measurement, not about hearing. A stage may be measurable and
inaudible, or---more troublingly---audible in ways no declared measurement
captures. The language's guarantees stop at the declared instruments.

\emph{Nothing here composes music.} Sangoma builds sounds against stated targets.
It does not choose the targets, and the choice is the interesting part.

%──────────────────────────────────────────────────────────────────────
\section{Conclusion}
\label{sec:conclusion}
%──────────────────────────────────────────────────────────────────────

We have specified a language in which a sound is declared by what it must
satisfy, at a stated resolution, rather than by the steps that produce it; in
which reachability is decided before rendering and unreachability is a diagnostic
carrying the depth required (\Cref{prop:reach}); in which a single stage cannot
be sharpened without bound, so that resolution is reached by composition
(\Cref{thm:refine}, \Cref{thm:compose}); and in which a stage engulfed by its own
separator is refused statically rather than rendered and found silent
(\Cref{thm:detect}).

The rules it enforces are physical---headroom, latency, resolution---and it is
deliberately silent above them. What it offers a producer is not judgement about
their chain but two facts they cannot presently obtain: whether a target is
reachable with the stages they have, and by how much they fall short if it is
not.

\bibliographystyle{plainnat}
\bibliography{references}

\end{document}
